%!TEX root = thesis.tex
\chapter{Introduction}
\label{chap:introduction}

% ============================================================
% PLACEHOLDER — text will be added section by section
% during the subchapter review pass with Gard.
% ============================================================

\section{Wind Turbines and the Importance of Operational Context}
\label{sec:intro:context}

Wind turbines are complex pieces of machinery that operate across a wide range of
conditions. At lower wind speeds, the turbine sits in standby not producing anything.
As the wind picks up, the control system gradually spins up the rotor, the generator
engages and the turbine begins producing power. At higher wind speeds, the turbine
operates in a partial load state where the goal is to extract as much energy as possible
from the wind. Above a rated wind speed limit, the turbine has reached its maximum power
output and the control system is actively limiting power production by adjusting the blade
pitch angle. At extreme wind speeds, the turbine shuts down entirely to protect itself
from damage.

Each of these operating states, and the transitions between them, produces its own
pattern in the SCADA data. Wind speed, rotor RPM and blade pitch angle all behave
according to different physical laws in each state. In partial load, for example, the
power output follows a cubic relationship with wind speed and blade pitch is held at
zero degrees to maximise aerodynamic efficiency. In full load, the same power output
is held constant while blade pitch increases actively to limit rotor torque. A model
that does not know which of these two states the turbine is in will struggle to make
sense of either one.

This is the core problem that motivates this thesis. Before any other analysis of
SCADA data can be done reliably, before fault detection, before power curve modelling,
before degradation tracking, the operating state of the turbine must be identified.
This thesis calls this problem operational zone classification.

\section{SCADA Systems in Wind Turbine Monitoring}
\label{sec:intro:scada}

% TODO

\section{Operational Zones: Definition and Importance}
\label{sec:intro:zones}

% TODO

\section{Prognostics, Health Management, and Digital Twins}
\label{sec:intro:phm}

% TODO

\section{Approaches to Zone Classification}
\label{sec:intro:approaches}

% TODO

\section{Research Gap and Motivation}
\label{sec:intro:gap}

% TODO

\section{Research Questions}
\label{sec:intro:rq}

% TODO

\section{Scope and Delimitations}
\label{sec:intro:scope}

% TODO

\section{Research Method}
\label{sec:intro:method}

% TODO

\section{Thesis Structure}
\label{sec:intro:structure}

% TODO
